\section{Background and Related Work}\label{sec:preliminary}

\paragraph{Generated search text.}
RAG combines a generator with an evidence store: a searchable corpus of
text that the system retrieves from before answering \citep{lewis2020rag}.
For legal tasks the goal is to surface the controlling authority--which are the legal statues/precedents needed to resolve a dispute--for a
given question, but the question and the corpus often use different
language. HyDE addresses this by generating a hypothetical document and
embedding that text for retrieval \citep{gao2023hyde}. Query2doc and related expansion methods
similarly use a language model to turn sparse query text into richer retrieval
handles \citep{wang2023query2doc}. For legal tasks, the useful generated text
is usually not a longer version of the question. It is a passage in the style
of the authority the retriever should find.

\paragraph{Legal RAG evaluation.}
Legal evaluation benchmarks and corpora have standardized many
legal-understanding tasks
\citep{guha2023legalbench,chalkidis2022lexglue,henderson2022pileoflaw}, and
recent work has emphasized the risk of unsupported legal assistant answers
\citep{magesh2024hallucinations}. Legal RAG benchmarks such as LegalBench-RAG
and Legal RAG Bench argue for evaluating retrieval and answering together
\citep{pipitone2024legalbenchrag,butler2026legalragbench}. We follow that joint
evaluation but split the failure analysis into two questions: did the
retriever miss the evidence, or did the answer model misuse it?
